% ===== 第二段：Agent能力与工具调用评测 =====
\subsubsection{基于大语言模型的智能体能力与工具调用评测研究现状}

LLM Agent研究的早期主线是如何使模型在推理过程中调用外部环境。
ReAct\cite{yao2023react}将推理轨迹与行动轨迹交替组织，提出"思考—行动—观察"（Thought-Action-Observation）循环范式，
使模型能够在与外部工具或检索系统的交互中完成具有现实侧效应的任务，奠定了工具调用型Agent的基本架构。
Toolformer\cite{schick2023toolformer}探索模型自监督学习何时调用API、调用何种API以及如何吸收工具返回结果，
证明仅需少量标注样本即可使语言模型掌握多种外部工具的使用，从而大幅降低了工具调用能力的训练成本。

Gorilla\cite{patil2024gorilla}面向大规模API调用构建了覆盖1600余个API的训练数据集和评测框架，
并提出通过检索增强解决API参数幻觉问题，说明工具调用能力已成为现代Agent系统的核心基础能力；
ToolLLM\cite{qin2024toolllm}则通过构建覆盖16000余个真实世界API的RapidAPI数据集，
系统评估了LLM在工具选择、参数生成和多步规划方面的综合能力。
CodeAct\cite{wang2024codeact}提出以可执行Python代码动作替代传统JSON格式的工具调用，
实验表明代码化动作能够显著减少多步交互轮次并提升任务完成率，为工具调用范式提供了新的设计维度。
Reflexion\cite{shinn2023reflexion}通过语言反馈机制使Agent能够从失败经验中学习，
将错误调用经验转化为自然语言反思并存入记忆，进一步提升了多步工具调用的鲁棒性与自适应能力。

在评测方面，AgentBench\cite{liu2024agentbench}将Agent能力放入操作系统、数据库、网络浏览等八个交互环境中进行系统评估，
揭示了当前LLM在多步工具调用任务上的普遍性能瓶颈。
WebArena\cite{zhou2023webarena}构建了高度真实的Web交互环境，强调任务的长时序性和目标多义性；
Mind2Web\cite{deng2023mind2web}则覆盖了超过2000个真实网页任务，
侧重在多样化网页布局和动态DOM结构下的泛化能力评测。
OSWorld\cite{xie2024osworld}与Windows Agent Arena\cite{bonatti2025windows}进一步把Agent放入真实操作系统环境中，
要求其执行文件操作、应用切换、GUI控制等复杂任务，验证了多步工具调用在非结构化现实场景下的挑战。
PEAR\cite{dong2026pear}等新近基准则将评测扩展至真实世界场景的多步执行，
覆盖更复杂的工具组合与跨应用操作，进一步揭示了Agent在安全敏感操作场景下的风险暴露面。

上述工作说明，工具型Agent的执行过程通常不是单次调用，而是包含多轮计划、调用、观察和修正的链式结构。
然而，现有能力评测基准普遍以任务完成率为核心指标，较少关注工具调用链的过程合规性——
即调用过程中是否存在权限越界、敏感数据泄露或不可逆操作等安全风险。
这一缺口意味着：一个任务完成率高的Agent系统，其工具调用链可能同时存在严重的安全隐患，
而现有评测框架无法有效识别和度量这类过程级风险。
